\section{问题重述}

\subsection{问题背景}

热风烘干通过调节烘房的温湿环境，使中药材经历预热平衡与恒温干燥两个阶段。
工艺参数的选择关系到干燥效率、能耗及成品品质，而依靠反复试验优化工艺需要较高的时间与成本投入。
因此，需要借助数学模型描述药材内部温度与水分浓度的时空变化，并据此确定满足烘干要求的处理时间。

题目所给药材近似呈圆柱形，初始长度为 $25\,\mathrm{cm}$、半径为 $2\,\mathrm{cm}$，
初始温度为 $28\,{}^\circ\mathrm{C}$，初始水分浓度为 $2.55\,\mathrm{kg/kg}$。
其中，药材的水分浓度指干基含水率。
附件1给出烘干初期烘房温度和水分浓度随时间的变化，附件2给出烘干过程中药材半径的变化，
题面附录2至附录4分别提供各问题所需的物性参数或经验公式。
在这些条件下，需要依次研究阶段内变化、全程变化、达标时间及尺寸收缩的影响。

\subsection{具体问题}

\textbf{问题一：预热平衡阶段的温度与水分浓度变化。}
根据附件1的烘房环境数据及题面附录2的参数，建立预热平衡阶段的数学模型，
求出前 $1800\,\mathrm{s}$ 内药材温度与水分浓度的时空分布。
在 $100$、$300$、$600$、$900$、$1200$、$1500$ 和 $1800\,\mathrm{s}$，
给出距药材中心 $0$、$0.5$、$1$、$1.5$ 和 $2\,\mathrm{cm}$ 处的两类结果。

\textbf{问题二：整个烘干过程的温度与水分浓度变化。}
将研究范围由预热平衡阶段扩展至包含恒温干燥阶段的全过程，
统一采用题面附录3的经验公式，描述密度、比热容、热传导系数随含水率的变化，
以及水分扩散系数对含水率和温度的依赖。
建立相应模型，并给出前 $3\,\mathrm{h}$ 内每隔 $0.5\,\mathrm{h}$、上述五个位置的温度与水分浓度。

\textbf{问题三：满足含水率要求的烘干时间。}
沿用问题二的全过程描述和题面附录3的经验公式，
以药材各处水分浓度均低于 $0.15\,\mathrm{kg/kg}$ 为结束条件，确定所需烘干时长，以小时计。
给出烘干过程中每隔 $6\,\mathrm{h}$ 及烘干结束时、距中心每隔 $0.5\,\mathrm{cm}$ 处的水分浓度。

\textbf{问题四：考虑尺寸收缩的烘干时间。}
在问题三的达标要求下，进一步考虑水分流失引起的药材尺寸变化，
依据附件2的半径数据及题面附录4的经验公式，确定药材收缩时的烘干时长。
给出每隔 $6\,\mathrm{h}$ 及烘干结束时的水分浓度，位置按距中心 $0.5\,\mathrm{cm}$ 的间隔选取，
并包含各时刻的药材表面。

题目要求的正文结果表及结果工作簿中的温度、含水率场值统一保留四位小数；连续事件时刻及数值检验指标按分析需要保留必要有效位数。正文结果按题面表1至表6的格式列示。
完整结果按附件3提供的 \texttt{result1.xlsx} 至 \texttt{result4.xlsx} 模板保存。
其中，问题一、二的完整结果采用 $1\,\mathrm{s}$ 的时间间隔，问题三、四采用 $60\,\mathrm{s}$ 的时间间隔，
空间间隔均为 $0.1\,\mathrm{cm}$。
